\subsection{Modeling Strategy}

The empirical analysis compares four models of increasing
flexibility. Each specification adds a distinct source of
predictive information or parameter variation:

\[
\text{Persistence}
\;\longrightarrow\;
\text{AR}
\;\longrightarrow\;
\text{Fixed HAR-X}
\;\longrightarrow\;
\text{TVP-HAR-X}.
\]

The persistence benchmark uses only current-hour volatility.
The autoregressive model estimates a flexible pattern of
short-run lag dependence. Fixed HAR-X replaces individual lags
with multi-horizon volatility summaries and supplements them
with market-activity variables. TVP-HAR-X retains the same
predictors while allowing their coefficients to evolve over
time.

This sequence separates three potential sources of forecast
improvement:

\begin{enumerate}[label=(\roman*)]
    \item persistence beyond the current hour;
    \item multi-horizon volatility and market-activity
    information;
    \item time variation in predictor coefficients.
\end{enumerate}

Table~\ref{tab:model-roles} summarizes the role of each model.

\begin{table}[H]
\centering
\caption{Forecasting models and their empirical roles}
\label{tab:model-roles}
\begin{tabular}{p{3.0cm}p{4.2cm}p{6.0cm}}
\toprule
Model
& Main information
& Empirical question \\
\midrule
Persistence
& Current-hour volatility
& Does the current volatility level provide an adequate
one-hour-ahead forecast? \\

AR(24)
& Twenty-four hourly lags
& Does a flexible short-run lag structure improve forecast
accuracy? \\

Fixed HAR-X
& Multi-horizon volatility and market-activity features
& Do structured long-horizon summaries and current market
conditions add predictive information? \\

TVP-HAR-X
& The HAR-X predictor set with evolving coefficients
& Does gradual coefficient variation improve forecasts relative
to a fixed relationship? \\
\bottomrule
\end{tabular}
\end{table}

\subsection{Persistence Benchmark}

The persistence forecast is

\begin{equation}
\widehat{y}_{t+1|t}^{\mathrm{Naive}}
=
y_t.
\label{eq:naive-model}
\end{equation}

On the original variance scale, this implies
\(\widehat{RV}_{t+1|t}^{\mathrm{Naive}}=RV_t\).

Volatility persistence makes this a relevant benchmark for
short-horizon forecasting. The model contains no estimated
parameters and represents the performance available from the
most recent observed volatility alone.

\subsection{Fixed-Parameter Autoregressive Model}

The general one-hour-ahead AR(\(p\)) model is

\begin{equation}
y_{t+1}
=
c
+
\sum_{j=1}^{p}
\phi_j y_{t-j+1}
+
\varepsilon_{t+1},
\label{eq:ar-model}
\end{equation}

where \(c\) is an intercept, \(\phi_j\) is the coefficient on
lag \(j\), and \(\varepsilon_{t+1}\) is a zero-mean forecast
disturbance.

Under this indexing convention,

\[
y_{t-j+1}
=
\begin{cases}
y_t, & j=1,\\
y_{t-1}, & j=2,\\
\vdots\\
y_{t-p+1}, & j=p.
\end{cases}
\]

The corresponding forecast is

\begin{equation}
\widehat{y}_{t+1|t}^{\mathrm{AR}(p)}
=
\widehat{c}
+
\sum_{j=1}^{p}
\widehat{\phi}_j y_{t-j+1}.
\label{eq:ar-forecast}
\end{equation}

The persistence benchmark is obtained from the restrictions
\(c=0\) and \(\phi_1=1\). Estimating the coefficients allows
the model to capture mean reversion and distribute predictive
weight across several recent hours.

\subsubsection{Stationarity Assessment}

The AR analysis begins with an assessment of the training-period
log realized-variance series. The graphical diagnostics include
the series itself, a 168-hour rolling mean, a 168-hour rolling
standard deviation, the autocorrelation function, and the
partial autocorrelation function.

The augmented Dickey--Fuller test accommodates higher-order
serial dependence through lagged differences
\citep{saiddickey1984}. Its null hypothesis is

\[
H_0^{\mathrm{ADF}}:
\text{the series contains a unit root}.
\]

The KPSS test uses level stationarity as the null hypothesis
\citep{kpss1992}:

\[
H_0^{\mathrm{KPSS}}:
\text{the series is level stationary}.
\]

The two tests provide complementary evidence. Joint rejection
can arise when a series is mean reverting but exhibits changes
in its level or dependence structure over long periods.

Formal tests are applied to the longest continuous valid
training segment. This avoids closing gaps by treating
nonadjacent observations as consecutive.

\subsubsection{Order Selection and Estimation}

Candidate orders are restricted to
\(p\in\{1,2,\ldots,24\}\). The upper bound covers one complete
day of hourly dependence while limiting the number of estimated
parameters.

Each candidate model is estimated by ordinary least squares on
the training sample. AIC and BIC are reported as in-sample
diagnostics, while validation QLIKE determines the selected
order. Validation RMSE and MAE serve as secondary criteria.

The procedure selects \(p=24\), yielding

\begin{equation}
y_{t+1}
=
c
+
\sum_{j=1}^{24}
\phi_j y_{t-j+1}
+
\varepsilon_{t+1}.
\label{eq:ar24-model}
\end{equation}

Dynamic stability is assessed through the roots of the
characteristic polynomial

\begin{equation}
1
-
\phi_1 z
-
\phi_2 z^2
-
\cdots
-
\phi_{24} z^{24}
=
0.
\label{eq:ar-characteristic-polynomial}
\end{equation}

The model is stable when all roots lie outside the unit circle.
Residual dependence is evaluated using autocorrelations and
Ljung--Box statistics at several horizons.

AR(24) serves as the principal fixed-lag benchmark for the
subsequent HAR-X comparisons.

\subsection{The HAR-X Specification}

AR(24) estimates a separate coefficient for each of the
previous 24 hourly observations and contains no direct summary
of volatility beyond one day. The HAR framework represents
persistent volatility through aggregated components measured
over several horizons \citep{corsi2009}.

The volatility predictor vector is

\[
\mathbf{h}_t
=
\begin{bmatrix}
HAR^{(1)}_t \\
HAR^{(6)}_t \\
HAR^{(24)}_t \\
HAR^{(168)}_t
\end{bmatrix}.
\]

The one-hour and six-hour components describe recent
conditions, while the 24-hour and 168-hour components summarize
daily and weekly volatility. This construction compresses the
lag history into four interpretable variables and avoids
estimating a separate coefficient for every hourly lag.

Three market-activity predictors are added:

\[
\mathbf{m}_t
=
\begin{bmatrix}
DS_t \\
Trades_t \\
Buy_t
\end{bmatrix}.
\]

Trading activity is motivated by the relationship between
information arrival, transaction intensity, and price
variability \citep{clark1973,jones1994}. Downside realized
semivariance measures variation generated by negative
intraday returns \citep{barndorffnielsen2010}. The
aggressive-buying proxy summarizes the direction of
taker-initiated volume.

These variables are derived from aggregated candlestick data.
They represent market-activity and order-flow proxies rather
than direct measures of bid--ask spreads, market depth, or
other limit-order-book conditions.

\subsection{Fixed HAR-X Model}

Let the standardized predictor vector be

\begin{equation}
\mathbf{x}_t
=
\begin{bmatrix}
1 \\
z_{HAR^{(1)},t} \\
z_{HAR^{(6)},t} \\
z_{HAR^{(24)},t} \\
z_{HAR^{(168)},t} \\
z_{DS,t} \\
z_{Trades,t} \\
z_{Buy,t}
\end{bmatrix}.
\label{eq:harx-design-vector}
\end{equation}

The first element is an intercept, and the remaining variables
are standardized using training-sample moments.

The fixed-parameter HAR-X model is

\begin{equation}
y_{t+1}
=
\mathbf{x}_t^{\mathsf{T}}
\boldsymbol{\beta}
+
\varepsilon_{t+1},
\qquad
\varepsilon_{t+1}
\sim
\mathcal{N}(0,H),
\label{eq:fixed-harx-vector}
\end{equation}

with coefficient vector

\[
\boldsymbol{\beta}
=
\begin{bmatrix}
\mu &
\beta_1 &
\beta_6 &
\beta_{24} &
\beta_{168} &
\beta_{DS} &
\beta_{Trades} &
\beta_{Buy}
\end{bmatrix}^{\mathsf{T}}.
\]

Equivalently,

\begin{align}
y_{t+1}
={}&
\mu
+
\beta_1 z_{HAR^{(1)},t}
+
\beta_6 z_{HAR^{(6)},t}
+
\beta_{24} z_{HAR^{(24)},t}
\nonumber\\
&+
\beta_{168} z_{HAR^{(168)},t}
+
\beta_{DS} z_{DS,t}
+
\beta_{Trades} z_{Trades,t}
+
\beta_{Buy} z_{Buy,t}
+
\varepsilon_{t+1}.
\label{eq:fixed-harx-expanded}
\end{align}

The coefficients are estimated by ordinary least squares.
Their interpretation is conditional on the remaining
predictors. The HAR terms measure incremental dependence at
different volatility horizons, while the three additional
variables capture information associated with the composition
and intensity of current market activity.

Coefficient inference uses Newey--West
heteroskedasticity- and autocorrelation-consistent standard
errors with 24 lags \citep{neweywest1987}. These standard
errors affect inference but leave the OLS coefficient estimates
and forecasts unchanged.

Fixed HAR-X provides the benchmark for evaluating the value of
time-varying coefficients under an identical observable
information set.

\subsection{Motivation for Time-Varying Parameters}

Fixed HAR-X imposes a constant coefficient vector,
\(\boldsymbol{\beta}_t=\boldsymbol{\beta}\), throughout the
sample. This restriction may be strong for Bitcoin markets,
where volatility conditions, participation, and trading
activity can change substantially. Previous evidence also
shows that jumps and structural breaks affect Bitcoin
volatility forecasts \citep{shen2020}.

A constant coefficient summarizes the average predictive
relationship over the estimation period. When that
relationship changes, a fixed model may adjust slowly and
combine observations generated under different market
conditions.

TVP-HAR-X treats the coefficient vector as a latent dynamic
state. Coefficients evolve gradually as new forecast errors
alter the estimated state distribution. This structure permits
adaptation while preserving continuity between adjacent
forecast origins.

\subsection{TVP-HAR-X State-Space Model}

TVP-HAR-X uses the same observable predictors as Fixed HAR-X
and allows their coefficients to vary over time. The model is
written in state-space form, linking the observed target to a
latent coefficient vector \citep{durbin2012}.

The observation equation is

\begin{equation}
y_{t+1}
=
\mathbf{x}_t^{\mathsf{T}}
\boldsymbol{\beta}_t
+
\varepsilon_{t+1},
\qquad
\varepsilon_{t+1}
\sim
\mathcal{N}(0,H).
\label{eq:tvp-observation}
\end{equation}

The state vector is

\begin{equation}
\boldsymbol{\beta}_t
=
\begin{bmatrix}
\mu_t \\
\beta_{1,t} \\
\beta_{6,t} \\
\beta_{24,t} \\
\beta_{168,t} \\
\beta_{DS,t} \\
\beta_{Trades,t} \\
\beta_{Buy,t}
\end{bmatrix}.
\label{eq:tvp-state-vector}
\end{equation}

Its transition equation is a random walk:

\begin{equation}
\boldsymbol{\beta}_{t+1}
=
\boldsymbol{\beta}_t
+
\boldsymbol{\eta}_{t+1},
\qquad
\boldsymbol{\eta}_{t+1}
\sim
\mathcal{N}
\left(
\mathbf{0},
\mathbf{Q}
\right).
\label{eq:tvp-state-transition}
\end{equation}

The state covariance matrix is restricted to

\begin{equation}
\mathbf{Q}
=
\operatorname{diag}
\left(
q_{\mu},
q_{\beta},
q_{\beta},
q_{\beta},
q_{\beta},
q_{\beta},
q_{\beta},
q_{\beta}
\right).
\label{eq:tvp-state-covariance}
\end{equation}

The intercept has innovation variance \(q_{\mu}\), and the
seven slope coefficients share \(q_{\beta}\). This restriction
reduces the variance parameters to \(H\), \(q_{\mu}\), and
\(q_{\beta}\), providing a parsimonious specification of
coefficient variation.

The state variances govern the rate of adaptation. Values close
to zero produce slowly changing coefficients, while larger
values permit faster responses to new observations. All three
variance parameters are estimated from the data.

\subsection{Kalman Filtering and Sequential Forecasting}

Conditional on the variance parameters, TVP-HAR-X is a linear
Gaussian state-space model. The Kalman filter recursively
produces forecasts and updates the conditional distribution of
the latent coefficients \citep{kalman1960,durbin2012}.

Let

\begin{equation}
\mathbf{a}_{t|t-1}
=
E
\left[
\boldsymbol{\beta}_t
\mid
\mathcal{F}_{t-1}
\right]
\label{eq:predicted-state}
\end{equation}

denote the predicted state, with covariance

\begin{equation}
\mathbf{P}_{t|t-1}
=
\operatorname{Var}
\left(
\boldsymbol{\beta}_t
\mid
\mathcal{F}_{t-1}
\right).
\label{eq:predicted-state-covariance}
\end{equation}

Under the random-walk transition,

\begin{equation}
\mathbf{a}_{t|t-1}
=
\mathbf{a}_{t-1|t-1},
\label{eq:kalman-state-prediction}
\end{equation}

and

\begin{equation}
\mathbf{P}_{t|t-1}
=
\mathbf{P}_{t-1|t-1}
+
\mathbf{Q}.
\label{eq:kalman-covariance-prediction}
\end{equation}

The one-hour-ahead forecast is

\begin{equation}
\widehat{y}_{t+1|t}
=
\mathbf{x}_t^{\mathsf{T}}
\mathbf{a}_{t|t-1},
\label{eq:kalman-observation-forecast}
\end{equation}

with conditional forecast variance

\begin{equation}
F_t
=
\mathbf{x}_t^{\mathsf{T}}
\mathbf{P}_{t|t-1}
\mathbf{x}_t
+
H.
\label{eq:kalman-forecast-variance}
\end{equation}

After \(y_{t+1}\) becomes available, the forecast innovation is

\begin{equation}
v_t
=
y_{t+1}
-
\widehat{y}_{t+1|t},
\label{eq:kalman-innovation}
\end{equation}

and the Kalman gain is

\begin{equation}
\mathbf{K}_t
=
\frac{
\mathbf{P}_{t|t-1}
\mathbf{x}_t
}{
F_t
}.
\label{eq:kalman-gain}
\end{equation}

The filtered state and covariance are updated according to

\begin{equation}
\mathbf{a}_{t|t}
=
\mathbf{a}_{t|t-1}
+
\mathbf{K}_t v_t,
\label{eq:kalman-state-update}
\end{equation}

and

\begin{equation}
\mathbf{P}_{t|t}
=
\mathbf{P}_{t|t-1}
-
\mathbf{K}_t
\mathbf{x}_t^{\mathsf{T}}
\mathbf{P}_{t|t-1}.
\label{eq:kalman-covariance-update}
\end{equation}

The magnitude of each update depends on the forecast
innovation, state uncertainty, predictor values, and
observation variance. The filter therefore regulates
coefficient adjustment through the Kalman gain.

\subsection{Missing Observations}

Invalid target hours remain on the hourly calendar. During
such an hour, the state prediction proceeds according to

\[
\mathbf{a}_{t|t-1}
=
\mathbf{a}_{t-1|t-1},
\qquad
\mathbf{P}_{t|t-1}
=
\mathbf{P}_{t-1|t-1}
+
\mathbf{Q}.
\]

Because no valid target is observed, the measurement update is
omitted:

\[
\mathbf{a}_{t|t}
=
\mathbf{a}_{t|t-1},
\qquad
\mathbf{P}_{t|t}
=
\mathbf{P}_{t|t-1}.
\]

This treatment preserves elapsed calendar time across missing
hours.

\subsection{Parameter Estimation}

The initial state mean is set to the Fixed HAR-X coefficient
estimate. The initial covariance is

\begin{equation}
\mathbf{P}_0
=
0.25\mathbf{I}_8,
\label{eq:initial-state-covariance}
\end{equation}

where \(\mathbf{I}_8\) is the \(8\times8\) identity matrix.

The variance-parameter vector is
\(\boldsymbol{\theta}=(H,q_{\mu},q_{\beta})^{\mathsf{T}}\).
It is estimated by maximizing the Kalman-filter
prediction-error likelihood \citep{durbin2012}.

Conditional on information available before observation \(t\),

\begin{equation}
v_t
\sim
\mathcal{N}(0,F_t).
\label{eq:innovation-distribution}
\end{equation}

Ignoring constants, the log likelihood is

\begin{equation}
\ell
\left(
\boldsymbol{\theta}
\right)
=
-\frac{1}{2}
\sum_{t \in \mathcal{T}_{\mathrm{valid}}}
\left[
\log(F_t)
+
\frac{v_t^2}{F_t}
\right],
\label{eq:kalman-log-likelihood}
\end{equation}

where \(\mathcal{T}_{\mathrm{valid}}\) contains observations
with valid targets.

Positivity is imposed through exponential transformations of
the variance parameters. Optimization is performed from
several initial values, and the converged solution with the
highest likelihood is retained.

During model development, the variance parameters are estimated
on the training sample and evaluated sequentially on
validation data. After the specification is fixed, they are
re-estimated using the combined training and validation
period for the final test forecasts.

\subsection{Interpretation of the Time-Varying Coefficients}

The estimated state paths describe conditional predictive
associations. Causal interpretation lies outside the scope of
the model.

For example, a positive value of \(\beta_{Trades,t}\) indicates
that higher trading activity is associated with higher
next-hour log realized variance at time \(t\), conditional on
the remaining predictors. Changes in \(\beta_{1,t}\) and
\(\beta_{168,t}\) describe variation in the relative
predictive importance of recent and weekly volatility.

The state estimates are used to assess coefficient stability,
identify predictors with substantial temporal variation, and
compare their behavior across market conditions.

\subsection{Model Diagnostics}

AR(24) diagnostics include characteristic roots, residual
autocorrelations, and Ljung--Box portmanteau statistics
\citep{ljungbox1978}.

For Fixed HAR-X, the analysis reports the residual mean,
variance, first-order autocorrelation, Durbin--Watson
statistic, and Ljung--Box statistics.

TVP-HAR-X diagnostics use standardized one-step-ahead
innovations,

\begin{equation}
u_t
=
\frac{v_t}{\sqrt{F_t}}.
\label{eq:standardized-innovation}
\end{equation}

Under the linear Gaussian specification, these innovations
should have approximately zero mean, unit variance, and limited
serial correlation.

Because the diagnostic samples are large, small residual
correlations may produce statistically significant
portmanteau tests. The test results are interpreted jointly
with residual magnitudes, short-lag autocorrelations, and
out-of-sample forecast performance.

\subsection{Final Model Comparison}

Model selection is conducted using the training and validation
periods. The selected AR(24), Fixed HAR-X, and TVP-HAR-X
specifications are then re-estimated on the combined
development sample and evaluated on the same 2025 test
timestamps.

The comparison addresses three questions:

\begin{enumerate}[label=(\arabic*)]
    \item Does AR(24) improve on volatility persistence?
    \item Does Fixed HAR-X improve on AR(24) through
    multi-horizon and market-activity information?
    \item Does coefficient variation improve on Fixed HAR-X
    under an identical predictor set?
\end{enumerate}

The final comparison isolates the contribution of parameter
variation because Fixed HAR-X and TVP-HAR-X use the same
observable information.